% year: תשע"ט
% semester: ב
% moed: ב
% number: 2א
% points: 15
% categories: ivt
% source: exams/מבחנים/תשעט/חודיסמן מועד ב תשעט.pdf

\begin{question}
הוכיחו כי למשוואה $\dfrac{1}{x+1}+\dfrac1x+\dfrac{1}{x-1}=1$ יש לפחות שני פתרונות.
\end{question}

\begin{hint}
הגדירו $g(x)=\frac{1}{x+1}+\frac1x+\frac1{x-1}-1$ ובדקו את הגבולות החד-צדדיים של $g$ בקצוות הקטעים $(-1,0)$ ו-$(0,1)$. השתמשו במשפט ערך הביניים.
\end{hint}

\begin{solution}
נגדיר $g(x)=\frac{1}{x+1}+\frac1x+\frac1{x-1}-1$, הרציפה בכל קטע שאינו מכיל את $-1,0,1$.

\textbf{הקטע $(-1,0)$:} כאשר $x\to-1^+$, המחובר $\frac1{x+1}\to+\infty$ ושאר המחוברים חסומים, ולכן $g(x)\to+\infty$. כאשר $x\to0^-$, $\frac1x\to-\infty$ ולכן $g(x)\to-\infty$. מכאן קיימות נקודות $-1<p<q<0$ עם $g(p)>0$ ו-$g(q)<0$. $g$ רציפה ב-$[p,q]$, ולפי משפט ערך הביניים קיימת $x_1\in(p,q)$ עם $g(x_1)=0$.

\textbf{הקטע $(0,1)$:} כאשר $x\to0^+$, $\frac1x\to+\infty$ ולכן $g\to+\infty$; כאשר $x\to1^-$, $\frac1{x-1}\to-\infty$ ולכן $g\to-\infty$. באותו אופן, לפי משפט ערך הביניים קיים $x_2\in(0,1)$ עם $g(x_2)=0$.

הקטעים זרים, ולכן $x_1\ne x_2$ ולמשוואה לפחות שני פתרונות. (באותו אופן יש פתרון שלישי ב-$(1,\infty)$: $g\to+\infty$ כאשר $x\to1^+$ ו-$g\to-1$ כאשר $x\to\infty$.)
\end{solution}
